\section{Introduction}
\label{sec:introduction}

Leading entropy records how fast a language grows, but it deliberately
forgets the bounded arithmetic left by an integer cutoff.  That loss is
especially visible for multiplicative shifts.  Their constraints do not run
along consecutive sites; they run along the chains
\[
  i,\;qi,\;q^2i,\ldots,\qquad q\nmid i.
\]
The standard chain product gives the leading term, yet a prefix ending at
\(N\) cuts the chains at different depths according to the base-\(q\)
digits of \(N\).  The natural question is therefore finite rather than
asymptotic: after subtracting the entropy term, which order-one values can
actually occur?

We answer this question exactly for every integer \(q\ge2\) and every finite
primitive zero--one adjacency matrix \(A\).  Let \(Z(N)\) be the number of
admissible labelings of \(\{1,\ldots,N\}\), let \(h\) be the leading prefix
entropy, and set \(E(N)=\log Z(N)-hN\).  The inverse limit
\[
  \mathbb Z_q:=\varprojlim_n \mathbb Z/q^n\mathbb Z
\]
is used as a compact state space even when \(q\) is composite; the boundary
map is real-valued, not \(q\)-adic-valued.  We prove that \(E(N)\) extends to
a continuous function \(E_{A,q}\colon\mathbb Z_q\to\mathbb R\) and that its
image is precisely the complete set of subsequential limits of \(E(N)\) as
\(N\to\infty\).

\begin{figure}[t]
  \centering
  \resizebox{\textwidth}{!}{\begin{tikzpicture}[
  >=Latex,
  every node/.style={font=\small},
  site/.style={circle,draw=chainblue,thick,minimum size=7mm,inner sep=1pt},
  newsite/.style={circle,draw=goldenorange,very thick,fill=goldenorange!12,minimum size=8mm,inner sep=1pt},
  box/.style={rounded corners,draw,thick,align=center,inner sep=5pt}
]
  \node[site] (i) at (0,2.0) {\(i\)};
  \node[site] (qi) at (1.5,2.0) {\(qi\)};
  \node (dots) at (2.9,2.0) {\(\cdots\)};
  \node[site] (prev) at (4.35,2.0) {\(q^{v-1}i\)};
  \node[newsite] (new) at (6.2,2.0) {\(q^vi=N\)};
  \draw[chainblue,thick,->] (i)--(qi);
  \draw[chainblue,thick,->] (qi)--(dots);
  \draw[chainblue,thick,->] (dots)--(prev);
  \draw[goldenorange,very thick,->] (prev)--(new);
  \node[align=center,font=\scriptsize] at (2.35,2.72)
    {one \(q\)-adic chain};
  \node[align=center,font=\scriptsize,text=goldenorange!80!black] at (6.2,2.72)
    {the only updated endpoint};

  \node[box,draw=chainblue] (increment) at (3.0,0.15)
    {\(\displaystyle \log Z(N)-\log Z(N-1)=c_{\nu_q(N)}\)};
  \draw[->,thick] (new.south) to[out=-90,in=15] (increment.east);

  \node[box,draw=boundaryteal,text width=5.5cm] (inverse) at (10.8,0.15)
    {\(\displaystyle (N\bmod q^v)_{v\ge1}\longrightarrow x\in\mathbb Z_q\)\\[1mm]
     \(\displaystyle E_{A,q}(x)=
       -\sum_{v\ge1}(d_v-d_{v-1})\frac{x\bmod q^v}{q^v}\)};
  \node[box,draw=proofgreen,text width=5.5cm] (image) at (10.8,-1.9)
    {\(\displaystyle \Acc_{N\to\infty}\{\log Z(N)-hN\}
       =E_{A,q}(\mathbb Z_q)\)};
  \draw[->,very thick,boundaryteal] (increment.east) --
    node[above,font=\scriptsize,fill=white,inner sep=1pt] {valuation sum}
    (inverse.west);
  \draw[->,very thick,proofgreen] (inverse.south) -- (image.north);
\end{tikzpicture}
}
  \caption{One cutoff step changes exactly one \(q\)-adic chain and hence
  exactly one valuation level.  Summing these levelwise changes gives a
  uniformly convergent residue series; continuity and explicit diverging
  representatives identify its real image with the complete finite-size
  accumulation set.}
  \label{fig:chain-boundary}
\end{figure}

The binary golden adjacency exposes additional geometry.  Its boundary
function is a signed digit series whose coefficients alternate and dominate
their entire future tails.  This all-level estimate makes the real boundary
image a strongly separated Cantor set and determines its Hausdorff and box
dimensions.  The same coefficient tails control an ordinary generating
function at primitive dyadic roots, producing a dense family of nonzero
radial leading coefficients and therefore a natural boundary.

\paragraph{Theorem contributions.}
The mathematical results are the following.
\begin{enumerate}[label=(T\arabic*)]
  \item An exact finite-\(N\) increment and residue identity, followed by the
  equality
  \[
    \Acc_{N\to\infty}\{\log Z(N)-hN\}=E_{A,q}(\mathbb Z_q).
  \]
  \item For the binary golden control, a full-tail separation theorem and
  \[
    \dim_H E_{A,2}(\mathbb Z_2)
    =\dim_B E_{A,2}(\mathbb Z_2)
    =\frac{\log2}{2\log\varphi}.
  \]
  \item For the ordinary cutoff series \(G(z)=\sum_{N\ge0}E(N)z^N\), the
  exact radial leading coefficient at every primitive dyadic root and the
  resulting unit-circle natural boundary.
\end{enumerate}

The multiplicative-shift object and its leading theory are established in
earlier work \citep{fan2012level,kenyon2012hausdorff,ban2019pattern}.
Accordingly, Fibonacci word counts, the chain product, leading entropy and
dimensions, and valid boundary-complexity results are background rather than
contributions of this paper.  \Cref{sec:prior} gives the direct source
comparison and reconciles a one-dimensional author-manuscript specialization
of \citet{ban2023boundary}, with an explicit version-of-record caveat.

The general boundary law, golden Cantor theorem, and natural-boundary theorem
appear in \cref{sec:exact,sec:golden,sec:natural}; scientific limitations are
collected in \cref{sec:scope}.  Two independent exact evaluators also replay
finite consequences of the formulas.  This reproducibility material is
placed in \cref{sec:replay}: agreement checks implementations only and is not
used to prove uniform convergence, the accumulation equality, separation, or
the natural boundary.  Artifact-specific Route and chronology records are
confined to \cref{app:route-chronology}.
